As a direct sensitivity test of the isotropic confinement scale, we repeated the continuum-certified calculation at $\omega/2\pi=1.10$, $1.20$, and $1.30~\mathrm{MHz}$, spanning the two experimental radial secular frequencies.  With the calibrated radius fixed at $r_c=25.876731~\mathrm{nm}$, the three negative-energy levels remain present throughout the interval.  The continuum ground-state energies are -15.065268, -15.000000, and -14.930063 MHz, respectively, corresponding to a total endpoint span of 0.135205 MHz.  The analogous spans of the first and second excited levels are 0.285445 and 0.498763 MHz.

To separate direct confinement sensitivity from the use of the single binding-energy calibration anchor, we also recalibrated $r_c$ independently at each frequency using the continuum condition $E_0/h=-15~\mathrm{MHz}$.  The required radii are 25.902273, 25.876731, and 25.849630 nm, a total variation of 0.052643 nm across the tested interval.  After recalibration, the negative-state count remains $N_-=3$, while the first excited level changes from -9.298517 to -9.119450 MHz and the second from -3.736071 to -3.313972 MHz, with endpoint spans of 0.179068 and 0.422099 MHz.  These shifts quantify sensitivity within the reduced stationary model and should not be interpreted as an uncertainty estimate for the full anisotropic Paul-trap dynamics.
